% =============================================================================
%  Capítulo 7 -- Conclusiones y trabajo futuro
% =============================================================================
\chapter{Conclusiones y trabajo futuro}
\label{cap:conclusiones}

\section{Conclusiones generales}
\label{sec:conclusiones-generales}

[Síntesis de los hallazgos principales de la tesis, organizada en torno a los
cuatro objetivos específicos planteados en la Sección~\ref{sec:objetivos}.]

\begin{enumerate}[label=\textbf{C\arabic*.},leftmargin=2.2em]
  \item \textbf{Respecto al objetivo O1} (documentación de clasismo
    estructural y discriminación lingüística en el contexto mexicano):
    [síntesis].
  \item \textbf{Respecto al objetivo O2} (diseño del \emph{benchmark} de
    textos controlados): [síntesis].
  \item \textbf{Respecto al objetivo O3} (evaluación experimental del
    comportamiento del sistema): [síntesis].
  \item \textbf{Respecto al objetivo O4} (diseño y cálculo del \INS):
    [síntesis].
\end{enumerate}

\section{Líneas de investigación futura}
\label{sec:trabajo-futuro}

\subsection{Validación empírica del \INS}

[Estudios sistemáticos para caracterizar el comportamiento del indicador en
distintas familias de modelos y \emph{benchmarks}, comparación con otras
métricas de equidad, calibración de umbrales orientativos.]

\subsection{Expansión del \textit{benchmark}}

[Incorporación de variantes regionales del español mexicano; cobertura de
tareas distintas a la clasificación binaria; ampliación a otros idiomas.]

\subsection{Mitigación de sesgos sociolectales}

[Aproximaciones de mitigación inspiradas en los hallazgos:
\textit{data augmentation} con paráfrasis controladas, entrenamiento con
restricciones de invariancia, ajustes de \textit{post-processing}.]

\subsection{Articulación con marcos regulatorios}

[Posibilidades de uso del \INS o de indicadores análogos en marcos
regulatorios emergentes sobre inteligencia artificial.]
